% ============================================================================
%  1. TỔNG QUAN
% ============================================================================
\section{Tổng quan}
\label{sec:tong-quan}

\subsection{Bối cảnh và động lực}

Luật Trật tự, an toàn giao thông đường bộ số 36/2024/QH15~\cite{luat36} và Nghị định
168/2024/NĐ-CP~\cite{nd168} cùng có hiệu lực từ ngày 01 tháng 01 năm 2025, thay thế toàn bộ
khung pháp lý cũ và nâng mạnh mức xử phạt ở nhiều nhóm hành vi. Chỉ trong hơn một năm sau
đó, khung này tiếp tục được sửa đổi hai lần: Luật 118/2025/QH15~\cite{luat118} (hiệu lực
01/7/2026) và Nghị định 238/2026/NĐ-CP~\cite{nd238} (hiệu lực 15/8/2026). Người tham gia giao
thông vì vậy đứng trước một khối văn bản vừa mới, vừa phân tán, vừa liên tục thay đổi.

Khó khăn thực tế không nằm ở chỗ thiếu thông tin mà ở chỗ thông tin bị \emph{tách rời}.
Luật quy định hành vi nào bị nghiêm cấm; nghị định mới quy định hành vi đó bị phạt bao
nhiêu tiền và bị trừ mấy điểm giấy phép lái xe. Câu hỏi thường gặp nhất của người dân ---
\emph{``vượt đèn đỏ phạt bao nhiêu''} --- do đó không thể trả lời nếu chỉ đọc một văn bản.
Các công cụ tìm kiếm phổ thông trả về toàn văn điều khoản, buộc người hỏi tự đối chiếu;
còn các trợ lý sinh ngôn ngữ tự nhiên thì có nguy cơ diễn giải sai lời của luật mà không
nêu được căn cứ.

Một câu trả lời sai trong tra cứu pháp luật không phải là ``kết quả kém hơn một chút'' mà
là thông tin dẫn người dân tới hành vi sai. Ba đặc điểm của lĩnh vực giao thông đường bộ
khiến nó đặc biệt phù hợp với hướng \textbf{tri thức tường minh}:
\begin{itemize}
  \item văn bản có cấu trúc chặt (điều -- khoản -- điểm) nên mô hình hoá được;
  \item tri thức mang tính \emph{suy diễn số} (ngưỡng nồng độ cồn, mức vượt tốc độ, điểm
        trừ cộng dồn) nên so khớp văn bản đơn thuần là không đủ;
  \item mọi câu trả lời đều phải truy nguyên được về một căn cứ pháp lý cụ thể.
\end{itemize}

\subsection{Phát biểu bài toán}
\label{sec:phat-bieu}

Cho cơ sở tri thức
\begin{equation}
  K = (C,\; R,\; \mathit{Rules},\; F,\; \mathit{Keyphrase})
  \label{eq:K}
\end{equation}
trong đó $C$ là tập khái niệm pháp lý, $R \subseteq C \times C$ là tập quan hệ hai ngôi,
$\mathit{Rules}$ là tập quy tắc giao thông, $F$ là tập sự kiện ``hành vi vi phạm -- chế
tài'', và $\mathit{Keyphrase}$ là từ điển cụm từ khoá ánh xạ ngôn ngữ người dùng vào định
danh tri thức. Mỗi phần tử tri thức $x$ mang một khối căn cứ
$cc(x) = (\text{văn bản}, \text{điều}, \text{khoản}, \text{điểm})$ và một khoảng hiệu lực
$hl(x) = [t_{\text{bắt đầu}},\, t_{\text{kết thúc}})$.

\begin{table}[H]
\centering
\caption{Định nghĩa hình thức đầu vào -- đầu ra của bài toán}
\label{tab:io}
\renewcommand{\arraystretch}{1.3}
\begin{tabularx}{\textwidth}{@{}>{\bfseries}p{2.7cm} L@{}}
\toprule
Thành phần & \textbf{Định nghĩa} \\
\midrule
Đầu vào & Truy vấn $q$ là chuỗi tiếng Việt tự nhiên (có dấu hoặc không dấu), kèm tuỳ chọn
mốc thời gian $t$ (mặc định: ngày hiện tại) và bộ lọc phương tiện. \\
Đầu ra & Bộ ba $A(q,t) = (p, T, E)$: $p \in \{P1,\dots,P7\}$ là lớp bài toán được nhận
diện; $T \subseteq C \cup \mathit{Rules} \cup F$ là danh sách tối đa $k$ mẩu tri thức đã
xếp hạng, mọi $x \in T$ đều thoả $t \in hl(x)$; $E$ là phần giải trình gồm nguyên văn điều
khoản và căn cứ $cc(x)$ của từng $x \in T$. \\
Ràng buộc & Hệ thống \emph{không} sinh ngôn ngữ tự nhiên: mọi câu chữ trong $E$ đều là
trích nguyên văn từ văn bản nguồn. Nếu không có $x$ nào thoả, hệ trả về tập rỗng kèm lý do
chứ không suy đoán. \\
Mục tiêu & Cực đại hoá xác suất đáp án chuẩn nằm ở hạng 1 của $T$ (Top-1) và nằm trong $T$
(Top-$k$), đồng thời giữ độ chính xác phân lớp $p$ và bảo đảm 100\% phần tử trả về truy
nguyên được căn cứ. \\
\bottomrule
\end{tabularx}
\end{table}

Phát biểu này tách bạch hai việc thường bị gộp làm một trong các hệ hỏi đáp: \textbf{tìm
đúng} mẩu tri thức ($T$) và \textbf{giải trình được} vì sao ($E$). Hệ thống bị ràng buộc
phải làm được cả hai, và chính ràng buộc ``không sinh ngôn ngữ'' là thứ phân biệt đồ án với
hướng RAG/LLM trình bày ở mục~\ref{sec:lien-quan}.

\subsection{Đóng góp của nhóm}

\begin{enumerate}
  \item \textbf{Cơ sở tri thức mở} cho lĩnh vực giao thông đường bộ theo khung pháp lý
        2024--2026: 73 khái niệm, 482 quan hệ, 109 quy tắc, 345 hành vi vi phạm và 1.674
        cụm từ khoá; 100\% mục có căn cứ pháp lý ở mức điều -- khoản -- điểm và có khoảng
        hiệu lực.
  \item \textbf{Mô hình hoá hiệu lực theo thời gian ở mức khoản}: sửa đổi được lưu như dữ
        liệu (63 bản ghi) chứ không hợp nhất cứng, nhờ đó trả lời được câu hỏi ``tại thời
        điểm $t$ thì áp dụng quy định nào''.
  \item \textbf{Thuật giải xử lý truy vấn sáu bước} với hàm điểm lai ba thành phần và bộ
        suy diễn số học cho các điều khoản phân khung theo ngưỡng; thí nghiệm loại bỏ thành
        phần chứng minh từng thành phần đều đóng góp thực chất.
  \item \textbf{Bộ dữ liệu đánh giá} 120 câu hỏi có đáp án chuẩn gán tới mức định danh tri
        thức, 40 truy vấn ngoài lĩnh vực, cùng toàn bộ kết quả đo, kiểm định thống kê và
        phân tích lỗi tái lập được bằng một lệnh chạy.
\end{enumerate}

\subsection{Cấu trúc báo cáo}

Mục~\ref{sec:lien-quan} khảo sát các công trình liên quan và xác định khoảng trống.
Mục~\ref{sec:du-lieu} trình bày cách xây dựng và phân tích bộ dữ liệu.
Mục~\ref{sec:phuong-phap} mô tả kiến trúc, thuật giải và thiết lập thực nghiệm.
Mục~\ref{sec:ket-qua} báo cáo kết quả, kiểm định thống kê và phân tích lỗi.
Mục~\ref{sec:ket-luan} trình bày ứng dụng, hạn chế, hướng phát triển và kết luận.
Phụ lục~A ghi phân công công việc; Phụ lục~B liệt kê các lệnh kiểm chứng lại số liệu.
